\SourceLeaf{527}of which a faint indication is given in the closing note of the Disquis. Arithm. p. 668 [\textsc{Gauss} Werke vol. I. p. 466].
\par This is a sufficient criterion for the case where $p$ has the form $8n+5$.
\par Let the number of classes formed by the binary forms in each of the two kinds for the determinant $-p$ be $=k$. The beginning of a table I constructed up to the determinant $-3000$ is in Disquis. Arithm. p. 520. [Art. 303.] It should also be noted that for a $p$ of the assumed form we always have $k=2m+1$ when $m$ denotes the number of decompositions of $p$ into three positive squares (I say positive to exclude $0$), as \textsc{Legendre} found by induction and as was first proved in the Disquis. Arithm. from the theory of \textit{ternary} forms. For example,
\begingroup\fontsize{8.7}{11}\selectfont\setlength{\arraycolsep}{2pt}
\[
\begin{array}{r*{10}{c}l}
\text{for }p=&5,&13,&29,&37,&53,&61,&101,&109,&149,&157&\text{etc.}\\
k=&1,&1,&3,&1,&3,&3,&7,&3,&7,&3&\\
m=&0,&0,&1,&0,&1,&1,&3,&1,&3,&1&\\[5pt]
&&&\overline{\begin{array}{c}4\\9\\16\end{array}}&\overline{\vphantom{\begin{array}{c}0\\0\\0\end{array}}\hspace{1em}}&\overline{\begin{array}{c}1\\16\\36\end{array}}&\overline{\begin{array}{c}9\\16\\36\end{array}}&\SourceVectorBrace{\begin{array}{ccc}1&4&16\\36&16&36\\64&81&49\end{array}}&\overline{\begin{array}{c}9\\36\\64\end{array}}&\SourceVectorBrace{\begin{array}{ccc}1&4&36\\4&64&64\\144&81&49\end{array}}&\overline{\begin{array}{c}4\\9\\144\end{array}}&
\end{array}
\]
\endgroup
With this assumption, the value of $b$ that is $\equiv\frac12rr\, (\mathrm{mod}.p)$ is always
\[
{}\equiv2k+a-1\equiv4m+a+1\ (\mathrm{mod}.8)
\]
which completely determines the sign of $b$. Here are 22 examples, doubling the extent of the table given at the end of the memoir.
\begin{center}
\begin{tabular}{r|r|r|r||r|r|r|r}
$p$&$k$&$a$&$b$&$p$&$k$&$a$&$b$\\\hline
$5$&$1$&$+1$&$+2$&$181$&$5$&$+9$&$+10$\\
$13$&$1$&$-3$&$-2$&$197$&$5$&$+1$&$-14$\\
$29$&$3$&$+5$&$+2$&$229$&$5$&$-15$&$+2$\\
$37$&$1$&$+1$&$-6$&$269$&$11$&$+13$&$+10$\\
$53$&$3$&$-7$&$-2$&$277$&$3$&$+9$&$+14$\\
$61$&$3$&$+5$&$-6$&$293$&$9$&$+17$&$+2$\\
$101$&$7$&$+1$&$-10$&$317$&$5$&$-11$&$+14$\\
$109$&$3$&$-3$&$+10$&$349$&$7$&$+5$&$+18$\\
$149$&$7$&$-7$&$-10$&$373$&$5$&$-7$&$+18$\\
$157$&$3$&$-11$&$-6$&$389$&$11$&$+17$&$-10$\\
$173$&$7$&$+13$&$+2$&$397$&$3$&$-19$&$-6$\\
\end{tabular}
\end{center}
